\documentclass[aspectratio=169,11pt]{beamer}

\usetheme{Madrid}
\usecolortheme{default}
\usefonttheme{professionalfonts}
\setbeamertemplate{navigation symbols}{}
\setbeamertemplate{footline}[frame number]

\usepackage{amsmath, amssymb, booktabs}

\title{Disability, Chronic Conditions, And Labor Supply}
\subtitle{Job Choice, Accommodation, And Dynamic Exit}
\author{Special Topic 6: Labor Market, Health, and Population -- Week 2}
\date{}

\begin{document}

\begin{frame}
  \titlepage
\end{frame}

\begin{frame}{Central question}
\begin{itemize}
    \item How do disability and chronic conditions reshape labor-force attachment, hours, occupation, accommodations, and exit?
    \item The labor objects are employment, hours, wages, occupation, job amenities, mobility, employer response, and welfare.
    \item Week 2 discipline: separate capacity, program rules, accommodation, local demand, and selection.
\end{itemize}
\end{frame}

\begin{frame}{Why this is labor economics}
\begin{itemize}
    \item Disability changes the feasible set for work, not only clinical status.
    \item Chronic conditions can alter hours, search, commuting, schedules, task assignment, and promotion.
    \item Firms respond through accommodation, screening, retention, reassignment, and job design.
    \item Program rules and benefit receipt matter because they change the employment boundary.
\end{itemize}
\end{frame}

\begin{frame}{Five mechanisms to keep separate}
\small
\begin{tabular}{p{0.22\linewidth} p{0.50\linewidth} p{0.18\linewidth}}
\toprule
Mechanism & What moves & Typical outcome \\
\midrule
Capacity & tasks, stamina, cognition, pain, commute & hours, exit \\
Programs & eligibility, benefits, reassessment risk & claiming, work \\
Accommodation & schedules, tasks, leave, remote work & retention \\
Local demand & feasible employers and occupations & mobility \\
Selection & who applies, receives, or is observed & bias \\
\bottomrule
\end{tabular}
\vspace{0.35cm}
\begin{itemize}
    \item A single employment coefficient can bundle all five mechanisms.
\end{itemize}
\end{frame}

\begin{frame}{Direct capacity effects vs program incentives}
\small
\[
V^W(H_{it}, A_{it}, D_{mt}) \geq V^B(H_{it}, R_{it}, B_{it})
\]
\begin{itemize}
    \item Capacity effects: health limits tasks, feasible hours, commuting, schedules, or productivity.
    \item Program effects: eligibility, earnings tests, waiting periods, benefits, and reassessment risks change the value of work.
    \item Disability status is not pure health; it blends impairment, rules, application costs, and local opportunities.
    \item Examiner, threshold, and reassessment designs identify program response for marginal cases.
\end{itemize}
\end{frame}

\begin{frame}{Dynamic constraints}
\small
\begin{tabular}{p{0.25\linewidth} p{0.32\linewidth} p{0.29\linewidth}}
\toprule
Object & Immediate margin & Dynamic margin \\
\midrule
Acute onset & absence, hours cut & recovery or persistent exit \\
Chronic condition & task and schedule friction & career scarring, downgrading \\
Functional limitation & feasible jobs shrink & lower mobility, amenity demand \\
Benefit receipt & award and work incentives & reentry and attachment \\
Accommodation & retention today & match-specific dependence \\
\bottomrule
\end{tabular}
\vspace{0.3cm}
\begin{itemize}
    \item The question is always: which horizon and which labor margin?
\end{itemize}
\end{frame}

\begin{frame}{Short-run vs long-run effects}
\begin{itemize}
    \item Short run: missed work, reduced hours, urgent medical spending, leave decisions, liquidity stress, and temporary accommodation.
    \item Medium run: job retention, employer reassignment, treatment response, benefit application, and search.
    \item Long run: persistent exit, occupational downgrading, lower income and consumption, weaker mobility, and welfare loss.
    \item A one-year employment estimate, five-year earnings estimate, and ten-year consumption estimate are different estimands.
\end{itemize}
\end{frame}

\begin{frame}{Treatment, accommodation, return to work}
\small
\begin{tabular}{p{0.24\linewidth} p{0.36\linewidth} p{0.26\linewidth}}
\toprule
Response & Labor-market channel & Main identification issue \\
\midrule
Treatment & improves capacity or stabilizes symptoms & selection into care \\
Accommodation & changes feasible tasks and schedules & endogenous firm response \\
Return-to-work policy & lowers reentry risk or benefit loss & program-margin LATE \\
Remote/flexible work & relaxes commute and schedule constraints & occupation sorting \\
\bottomrule
\end{tabular}
\vspace{0.35cm}
\begin{itemize}
    \item These are not outside the labor model; they change the feasible set and the value of work.
\end{itemize}
\end{frame}

\begin{frame}{Employer response and local demand}
\begin{itemize}
    \item Strong local demand can absorb health limitations through sorting, flexibility, and outside offers.
    \item Weak local demand can turn the same limitation into nonemployment or benefit application.
    \item Firms may retain and accommodate valuable workers, or screen and assign health-limited workers to lower-return tasks.
    \item Selection into observed disability status depends on severity, job loss, information, application costs, and expected benefits.
\end{itemize}
\end{frame}

\begin{frame}{Monopsony, job lock, and worker inelasticity}
\small
\begin{tabular}{p{0.28\linewidth} p{0.38\linewidth} p{0.20\linewidth}}
\toprule
Channel & Why outside options shrink & Labor implication \\
\midrule
Insurance lock & leaving risks coverage and treatment continuity & lower quits \\
Accommodation lock & current setup may not transfer & employer-specific rents \\
Local scarcity & few feasible jobs or commutes & lower elasticity \\
Stigma/uncertainty & new employers may not accommodate & search costs \\
Information advantage & current firm knows productivity and limits & wage-setting power \\
\bottomrule
\end{tabular}
\vspace{0.25cm}
\begin{itemize}
    \item Lower mobility may reveal lock-in, not better matches.
\end{itemize}
\end{frame}

\begin{frame}{Frontier methods}
\small
\begin{tabular}{p{0.30\linewidth} p{0.46\linewidth} p{0.14\linewidth}}
\toprule
Design & What it can identify & Caution \\
\midrule
Acute-shock event study & disruption and recovery after sharp health events & severity \\
Disability-onset admin event study & dynamic paths after observed onset & institutional status \\
Threshold/rule design & program incentives and receipt effects & local margin \\
Twin/sibling comparison & family endowment confounding & external validity \\
Genetic/MR design & chronic-condition channels & exclusion restriction \\
Matched employer-employee & firm accommodation and retention heterogeneity & sorting \\
Treatment reform & care, accommodation, or reentry response & bundled changes \\
\bottomrule
\end{tabular}
\end{frame}

\begin{frame}{Research Lab design}
\begin{itemize}
    \item Primary anchor: Meyer and Mok on long-run disability incidence for earnings, income, and consumption.
    \item Challenge anchor: Hill, Maestas, and Mullen on employer accommodation and labor supply.
    \item Reproduce: estimate a synthetic disability-onset event-study profile.
    \item Diagnose: classify direct capacity, program, accommodation, local-demand, and selection channels.
    \item Transfer: redesign the logic for acute shocks, treatment access, program thresholds, or matched-firm accommodation.
\end{itemize}
\end{frame}

\begin{frame}{Research design checklist}
\begin{enumerate}
    \item Name the health object: onset, diagnosis, function, chronic burden, certification, or benefit receipt.
    \item Name the labor margin: employment, hours, earnings, occupation, mobility, accommodation, or welfare.
    \item State the horizon: immediate disruption, medium-run adjustment, or long-run scarring.
    \item Separate capacity from programs, firms, local demand, and selection.
    \item Say what welfare object remains invisible in wages or employment.
\end{enumerate}
\end{frame}

\begin{frame}{Bridge to Week 3}
\begin{itemize}
    \item Week 2 centers working-age health limitations and the worker-employer match.
    \item Week 3 broadens the unit to household timing and population events.
    \item The next objects are fertility, mortality, caregiving, aging, and lifecycle labor allocation.
\end{itemize}
\end{frame}

\begin{frame}{Takeaways}
\begin{itemize}
    \item Disability and chronic conditions change capacity, job choice, amenities, program interaction, and employer response.
    \item Short-run disruption and long-run scarring are different estimands.
    \item Accommodation and treatment can change labor supply; they are not merely background policy.
    \item Job lock and accommodation dependence can make labor supply more inelastic and strengthen employer power.
\end{itemize}
\end{frame}

\end{document}
